
\documentclass[10pt, twocolumn]{article}
\usepackage{graphicx}
\usepackage{amsmath}
\usepackage{amssymb}
\usepackage{amsfonts}
\usepackage{textcomp}
\usepackage{amsthm}

%% ---- Theorem Environments ---------------------------------
\theoremstyle{plain}
\newtheorem{theorem}{Theorem}[section]
\newtheorem{lemma}[theorem]{Lemma}
\newtheorem{corollary}[theorem]{Corollary}
\newtheorem{proposition}[theorem]{Proposition}

\theoremstyle{definition}
\newtheorem{definition}[theorem]{Definition}
\newtheorem{example}[theorem]{Example}

\theoremstyle{remark}
\newtheorem{remark}[theorem]{Remark}

\title{P-02}
\author{A. E. Mahrouss\\amlal@nekernel.org}
\date{\today}

\begin{document}
\maketitle

\section{Abstract}

The document illustrates two theorems written down for note-writing purposes.

\section{Introduction and Concepts}


\section{Theorems and Proofs}
\begin{theorem}
There exists $\forall x \in \mathbb{P}$:
\begin{equation}
    \mathbb{P}_0 < \mathbb{P}_{n-1} \xrightarrow{} (x) < (x+1)
\end{equation}
s.t. $\forall x \in \mathbb{P}$:
\begin{equation}
    \{ x, \space \cdots, \space (x+n) \} \in \mathbb{P}
\end{equation}
\end{theorem}

\nocite{*}

\bibliographystyle{acm}
\bibliography{book1}

\end{document}